<% extends "base_template.tex" %>

<% block body %>

\thispagestyle{plain}
\begin{center}
  {\sffamily\bfseries\huge\color{WgsBlue} 肠道菌群健康评估报告} \\[0.6em]
  {\color{graytext}\small 样本编号：<{ sample_id }> \quad|\quad 生成时间：<{ report_date }>}
\end{center}

\begin{resultbox}{微生态综合健康评分}
  {\sffamily\bfseries\fontsize{34pt}{38pt}\selectfont\color{WgsBlue} <{ overall_score }>} \, {\large\color{graytext}/ 100} \\[0.5em]
  {\bfseries\color{darkgray} <{ rating }>}
\end{resultbox}

\begin{infobox}{核心摘要}
已检测物种数：\textbf{<{ n_total }> }；正常：\textbf{<{ n_normal }> }；偏高：\textbf{<{ n_high }> }；偏低：\textbf{<{ n_low }> }。
菌群结构评分为 \textbf{<{ dim_species }>}，疾病风险维度评分为 \textbf{<{ dim_disease }>}，药物交互维度评分为 \textbf{<{ dim_drug }>}。
\end{infobox}

\clearpage
\tableofcontents
\clearpage

<% if enterotype %>
\section{肠型分类分析}
\begin{infobox}{肠型定位}
当前肠型：\textbf{<{ enterotype.enterotype }>}；置信度：\textbf{<{ enterotype.confidence }>\%}。
<% if enterotype.description %>

<{ enterotype.description }>
<% endif %>
<% if enterotype_silhouette %>

肠型聚类轮廓系数（Silhouette Score）：\textbf{<{ enterotype_silhouette }> }。
<% endif %>
\end{infobox}

<% if enterotype_distribution %>
\subsection{肠型分布}
{\footnotesize
\begin{tabularx}{\textwidth}{L{0.38\textwidth} C{0.18\textwidth} Y}
\toprule
\rowcolor{WgsBlue!12}
\textbf{肠型} & \textbf{样本数} & \textbf{占比} \\
\midrule
<% for row in enterotype_distribution %>
<{ row.label }> & <{ row.count }> & <{ row.percent }> \\
<% if not loop.last %>\rowcolor{rowgray}<% endif %>
<% endfor %>
\bottomrule
\end{tabularx}
}
<% endif %>

<% if enterotype_marker_basis or enterotype_top_genus %>
\subsection{分类依据}
<% if enterotype_marker_basis %>
{\TableFont
\begin{tabularx}{\textwidth}{L{0.58\textwidth} C{0.24\textwidth}}
\toprule
\rowcolor{WgsBlue!12}
\textbf{标志性菌属组合} & \textbf{相对占比(\%)} \\
\midrule
<% for row in enterotype_marker_basis %>
<{ row.marker }> & <{ row.ratio }> \\
<% if not loop.last %>\rowcolor{rowgray}<% endif %>
<% endfor %>
\bottomrule
\end{tabularx}
}
<% endif %>

<% if enterotype_top_genus %>
{\TableFont
\begin{tabularx}{\textwidth}{L{0.58\textwidth} C{0.24\textwidth}}
\toprule
\rowcolor{WgsBlue!12}
\textbf{丰度靠前菌属} & \textbf{丰度(\%)} \\
\midrule
<% for row in enterotype_top_genus %>
<{ row.genus }> & <{ row.abundance }> \\
<% if not loop.last %>\rowcolor{rowgray}<% endif %>
<% endfor %>
\bottomrule
\end{tabularx}
}
<% endif %>
<% endif %>

<% if enterotype_cluster_plot %>
\subsection{t-SNE 聚类可视化}
\begin{figure}[H]
  \centering
  \includegraphics[width=0.88\textwidth]{\detokenize{<{ enterotype_cluster_plot }>}}
  \caption{基于菌属丰度的肠型 t-SNE 聚类图，当前样本已单独高亮标注}
\end{figure}
<% if enterotype_plot_interpretation %>
\begin{infobox}{图形解读}
<{ enterotype_plot_interpretation }>
\end{infobox}
<% endif %>
<% endif %>

<% if enterotype_relation_summary %>
\begin{recommendbox}{肠型结果解读}
<{ enterotype_relation_summary }>
\end{recommendbox}
<% endif %>
<% if enterotype_patient_interpretation %>
\begin{infobox}{患者解读}
<{ enterotype_patient_interpretation }>
\end{infobox}
<% endif %>
<% endif %>

<% if abnormal_species or normal_top_species or beneficial_species or harmful_species %>
\section{肠道菌群组成分析}
以下内容按有益菌与有害菌分层展示丰度、参考范围与说明；所有表格均使用固定列宽和自动换行，避免文字重叠。

<% if abnormal_species %>
\subsection{异常物种详情}
{\DenseTableFont
\begin{xltabular}{\textwidth}{P{0.24\textwidth} P{0.12\textwidth} C{0.10\textwidth} C{0.10\textwidth} C{0.10\textwidth} Z}
\toprule
\rowcolor{WgsBlue!12}
\textbf{物种名称} & \textbf{中文名} & \textbf{丰度(\%)} & \textbf{状态} & \textbf{偏差} & \textbf{健康解读} \\
\midrule
\endfirsthead
\rowcolor{WgsBlue!12}
\textbf{物种名称} & \textbf{中文名} & \textbf{丰度(\%)} & \textbf{状态} & \textbf{偏差} & \textbf{健康解读} \\
\midrule
\endhead
\midrule
\multicolumn{6}{r}{\footnotesize\color{graytext} 续下页}\\
\endfoot
\bottomrule
\endlastfoot
<% for row in abnormal_species %>
<{ row.species }> & <{ row.chinese_name }> & <{ row.abundance }> & <{ row.status }> & <{ row.deviation }> & <{ row.summary }> \\
<% if not loop.last %>\rowcolor{rowgray}<% endif %>
<% endfor %>
\end{xltabular}
}
<% endif %>

<% if normal_top_species %>
\subsection{正常范围内 Top 10 物种}
{\TableFont
\begin{xltabular}{\textwidth}{P{0.36\textwidth} P{0.18\textwidth} C{0.14\textwidth} Z}
\toprule
\rowcolor{WgsBlue!12}
\textbf{物种名称} & \textbf{中文名} & \textbf{丰度(\%)} & \textbf{参考范围 P5-P95} \\
\midrule
\endfirsthead
\rowcolor{WgsBlue!12}
\textbf{物种名称} & \textbf{中文名} & \textbf{丰度(\%)} & \textbf{参考范围 P5-P95} \\
\midrule
\endhead
\midrule
\multicolumn{4}{r}{\footnotesize\color{graytext} 续下页}\\
\endfoot
\bottomrule
\endlastfoot
<% for row in normal_top_species %>
<{ row.species }> & <{ row.chinese_name }> & <{ row.abundance }> & <{ row.ref_range }> \\
<% if not loop.last %>\rowcolor{rowgray}<% endif %>
<% endfor %>
\end{xltabular}
}
<% endif %>

<% if beneficial_species %>
\subsection{有益菌}
{\TableFont
\begin{xltabular}{\textwidth}{P{0.26\textwidth} P{0.15\textwidth} C{0.12\textwidth} C{0.17\textwidth} Z}
\toprule
\rowcolor{WgsBlue!12}
\textbf{物种名称} & \textbf{中文名} & \textbf{丰度(\%)} & \textbf{参考范围} & \textbf{说明} \\
\midrule
\endfirsthead
\rowcolor{WgsBlue!12}
\textbf{物种名称} & \textbf{中文名} & \textbf{丰度(\%)} & \textbf{参考范围} & \textbf{说明} \\
\midrule
\endhead
\midrule
\multicolumn{5}{r}{\footnotesize\color{graytext} 续下页}\\
\endfoot
\bottomrule
\endlastfoot
<% for row in beneficial_species %>
<{ row.species }> & <{ row.chinese_name }> & <{ row.abundance }> & <{ row.ref_range }> & <{ row.summary }> \\
<% if not loop.last %>\rowcolor{rowgray}<% endif %>
<% endfor %>
\end{xltabular}
}
<% endif %>

<% if harmful_species %>
\subsection{有害菌}
{\TableFont
\begin{xltabular}{\textwidth}{P{0.26\textwidth} P{0.15\textwidth} C{0.12\textwidth} C{0.17\textwidth} Z}
\toprule
\rowcolor{WgsBlue!12}
\textbf{物种名称} & \textbf{中文名} & \textbf{丰度(\%)} & \textbf{参考范围} & \textbf{说明} \\
\midrule
\endfirsthead
\rowcolor{WgsBlue!12}
\textbf{物种名称} & \textbf{中文名} & \textbf{丰度(\%)} & \textbf{参考范围} & \textbf{说明} \\
\midrule
\endhead
\midrule
\multicolumn{5}{r}{\footnotesize\color{graytext} 续下页}\\
\endfoot
\bottomrule
\endlastfoot
<% for row in harmful_species %>
<{ row.species }> & <{ row.chinese_name }> & <{ row.abundance }> & <{ row.ref_range }> & <{ row.summary }> \\
<% if not loop.last %>\rowcolor{rowgray}<% endif %>
<% endfor %>
\end{xltabular}
}
<% endif %>
<% if species_patient_interpretation %>
\begin{infobox}{患者解读}
<{ species_patient_interpretation }>
\end{infobox}
<% endif %>
<% endif %>

<% if pathway_data %>
\section{代谢通路活性分析}
共检测到 \textbf{<{ pathway_count }>} 条代谢通路，以下展示重点通路：

{\TableFont
\begin{xltabular}{\textwidth}{P{0.42\textwidth} C{0.14\textwidth} Z}
\toprule
\rowcolor{WgsBlue!12}
\textbf{通路名称} & \textbf{丰度} & \textbf{生物学提示} \\
\midrule
\endfirsthead
\rowcolor{WgsBlue!12}
\textbf{通路名称} & \textbf{丰度} & \textbf{生物学提示} \\
\midrule
\endhead
\midrule
\multicolumn{3}{r}{\footnotesize\color{graytext} 续下页}\\
\endfoot
\bottomrule
\endlastfoot
<% for row in pathway_data %>
<{ row.name }> & <{ row.abundance }> & <{ row.implication }> \\
<% if not loop.last %>\rowcolor{rowgray}<% endif %>
<% endfor %>
\end{xltabular}
}
<% if pathway_patient_interpretation %>
\begin{infobox}{患者解读}
<{ pathway_patient_interpretation }>
\end{infobox}
<% endif %>
<% endif %>

<% if metabolite_data %>
\section{代谢物含量分析}
共检测到 \textbf{<{ metabolite_count }>} 项代谢物，以下列出重点项目：

{\TableFont
\begin{xltabular}{\textwidth}{P{0.43\textwidth} C{0.16\textwidth} Z}
\toprule
\rowcolor{WgsBlue!12}
\textbf{代谢物} & \textbf{丰度} & \textbf{参考范围} \\
\midrule
\endfirsthead
\rowcolor{WgsBlue!12}
\textbf{代谢物} & \textbf{丰度} & \textbf{参考范围} \\
\midrule
\endhead
\midrule
\multicolumn{3}{r}{\footnotesize\color{graytext} 续下页}\\
\endfoot
\bottomrule
\endlastfoot
<% for row in metabolite_data %>
<{ row.name }> & <{ row.abundance }> & <{ row.ref_range }> \\
<% if not loop.last %>\rowcolor{rowgray}<% endif %>
<% endfor %>
\end{xltabular}
}
<% if metabolite_patient_interpretation %>
\begin{infobox}{患者解读}
<{ metabolite_patient_interpretation }>
\end{infobox}
<% endif %>
<% endif %>

<% if drug_data %>
\section{药物-微生物组交互分析}
共分析 \textbf{<{ drug_count }>} 项药物交互，较高风险项 \textbf{<{ high_risk_count }>} 个。

<% if high_risk_drugs %>
\subsection{高风险药物}
{\TableFont
\begin{tabularx}{\textwidth}{P{0.74\textwidth} C{0.18\textwidth}}
\toprule
\rowcolor{WgsBlue!12}
\textbf{药物名称} & \textbf{风险评分} \\
\midrule
<% for row in high_risk_drugs %>
<{ row.drug }> & <{ row.score }> \\
<% if not loop.last %>\rowcolor{rowgray}<% endif %>
<% endfor %>
\bottomrule
\end{tabularx}
}
<% endif %>

<% if all_drugs %>
\subsection{药物明细}
{\DenseTableFont
\begin{xltabular}{\textwidth}{P{0.60\textwidth} C{0.13\textwidth} Z}
\toprule
\rowcolor{WgsBlue!12}
\textbf{药物名称} & \textbf{评分} & \textbf{风险等级} \\
\midrule
\endfirsthead
\rowcolor{WgsBlue!12}
\textbf{药物名称} & \textbf{评分} & \textbf{风险等级} \\
\midrule
\endhead
\midrule
\multicolumn{3}{r}{\footnotesize\color{graytext} 续下页}\\
\endfoot
\bottomrule
\endlastfoot
<% for row in all_drugs %>
<{ row.drug }> & <{ row.score }> & <{ row.risk_level }> \\
<% if not loop.last %>\rowcolor{rowgray}<% endif %>
<% endfor %>
\end{xltabular}
}
<% endif %>

<% if drug_advice %>
\begin{recommendbox}{用药建议}
\begin{itemize}
<% for item in drug_advice %>
  \item <{ item }>
<% endfor %>
\end{itemize}
\end{recommendbox}
<% endif %>
<% if drug_patient_interpretation %>
\begin{infobox}{患者解读}
<{ drug_patient_interpretation }>
\end{infobox}
<% endif %>
<% endif %>

<% if high_risk_diseases or mid_risk_diseases or low_risk_diseases or disease_empty_message %>
\section{肠道菌群-疾病关联分析}

<% if high_risk_diseases %>
\subsection{高风险疾病}
{\DenseTableFont
\begin{xltabular}{\textwidth}{P{0.24\textwidth} C{0.11\textwidth} C{0.12\textwidth} C{0.12\textwidth} Z}
\toprule
\rowcolor{WgsBlue!12}
\textbf{疾病} & \textbf{评分} & \textbf{匹配物种} & \textbf{方向一致} & \textbf{关键物种} \\
\midrule
\endfirsthead
\rowcolor{WgsBlue!12}
\textbf{疾病} & \textbf{评分} & \textbf{匹配物种} & \textbf{方向一致} & \textbf{关键物种} \\
\midrule
\endhead
\midrule
\multicolumn{5}{r}{\footnotesize\color{graytext} 续下页}\\
\endfoot
\bottomrule
\endlastfoot
<% for row in high_risk_diseases %>
<{ row.disease }> & <{ row.risk_score }> & <{ row.matched_species }> & <{ row.agree_count }> & <{ row.key_species }> \\
<% if not loop.last %>\rowcolor{rowgray}<% endif %>
<% endfor %>
\end{xltabular}
}
<% endif %>

<% if mid_risk_diseases %>
\subsection{中风险疾病}
{\TableFont
\begin{tabularx}{\textwidth}{P{0.45\textwidth} C{0.12\textwidth} C{0.17\textwidth} Y}
\toprule
\rowcolor{WgsBlue!12}
\textbf{疾病} & \textbf{评分} & \textbf{匹配物种} & \textbf{方向一致} \\
\midrule
<% for row in mid_risk_diseases %>
<{ row.disease }> & <{ row.risk_score }> & <{ row.matched_species }> & <{ row.agree_count }> \\
<% if not loop.last %>\rowcolor{rowgray}<% endif %>
<% endfor %>
\bottomrule
\end{tabularx}
}
<% endif %>

<% if low_risk_diseases %>
\subsection{低风险疾病}
{\TableFont
\begin{tabularx}{\textwidth}{P{0.45\textwidth} C{0.12\textwidth} C{0.17\textwidth} Y}
\toprule
\rowcolor{WgsBlue!12}
\textbf{疾病} & \textbf{评分} & \textbf{匹配物种} & \textbf{方向一致} \\
\midrule
<% for row in low_risk_diseases %>
<{ row.disease }> & <{ row.risk_score }> & <{ row.matched_species }> & <{ row.agree_count }> \\
<% if not loop.last %>\rowcolor{rowgray}<% endif %>
<% endfor %>
\bottomrule
\end{tabularx}
}
<% endif %>

<% if disease_empty_message %>
\begin{infobox}{疾病关联状态}
<{ disease_empty_message }>
\end{infobox}
<% endif %>

\begin{infobox}{风险评分说明}
<{ disease_risk_note }>

<{ disease_disclaimer }>
\end{infobox}
<% if disease_patient_interpretation %>
\begin{infobox}{患者解读}
<{ disease_patient_interpretation }>
\end{infobox}
<% endif %>
<% endif %>

\section{综合健康评估与建议}
\begin{resultbox}{综合健康评分}
  {\sffamily\bfseries\fontsize{26pt}{30pt}\selectfont\color{WgsBlue} <{ overall_score }>} \, {\large\color{graytext}/ 100} \\[0.4em]
  {\bfseries\color{darkgray} <{ rating }>}
\end{resultbox}

{\footnotesize
\begin{tabularx}{\textwidth}{L{0.26\textwidth} C{0.14\textwidth} Y}
\toprule
\rowcolor{WgsBlue!12}
\textbf{维度} & \textbf{评分} & \textbf{说明} \\
\midrule
物种健康度 & <{ dim_species }> & 健康参考范围内物种占比越高，得分越高。 \\
\rowcolor{rowgray}
菌群稳定性 & <{ dim_stability }> & 异常物种越少，菌群结构越稳定。 \\
疾病防御力 & <{ dim_disease }> & 基于疾病关联风险的综合扣分结果。 \\
\rowcolor{rowgray}
药物耐受度 & <{ dim_drug }> & 药物与菌群潜在交互干扰越低，得分越高。 \\
\bottomrule
\end{tabularx}
}
<% if overall_patient_interpretation %>
\begin{infobox}{患者解读}
<{ overall_patient_interpretation }>
\end{infobox}
<% endif %>

<% if fmt_ok or fmt_degraded or fmt_unavailable %>
\subsection{FMT 配型预测}
<% if fmt_ok %>
\begin{recommendbox}{FMT 候选供体}
<{ fmt_message }>

<{ fmt_best_comment }>
\end{recommendbox}

<% if fmt_matches %>
{\TableFont
\begin{xltabular}{\textwidth}{C{0.07\textwidth} P{0.22\textwidth} C{0.13\textwidth} C{0.13\textwidth} C{0.11\textwidth} Z}
\toprule
\rowcolor{WgsBlue!12}
\textbf{排名} & \textbf{供体 ID} & \textbf{相对分} & \textbf{距离} & \textbf{病例} & \textbf{响应率} \\
\midrule
\endfirsthead
\rowcolor{WgsBlue!12}
\textbf{排名} & \textbf{供体 ID} & \textbf{相对分} & \textbf{距离} & \textbf{病例} & \textbf{响应率} \\
\midrule
\endhead
\midrule
\multicolumn{6}{r}{\footnotesize\color{graytext} 续下页}\\
\endfoot
\bottomrule
\endlastfoot
<% for row in fmt_matches %>
<{ row.rank }> & <{ row.donor_id }> & <{ row.score }> & <{ row.distance }> & <{ row.cases }> & <{ row.response_rate }> \\
<% if not loop.last %>\rowcolor{rowgray}<% endif %>
<% endfor %>
\end{xltabular}
}
<% endif %>
<% if fmt_quality_note %>
\begin{infobox}{模型输入质量}
<{ fmt_quality_note }>
\end{infobox}
<% endif %>
<% if fmt_tsne_plot %>
\subsection{FMT 配型嵌入可视化}
\begin{figure}[H]
  \centering
  \includegraphics[width=0.88\textwidth]{\detokenize{<{ fmt_tsne_plot }>}}
  \caption{基于三元组模型嵌入空间的 FMT 候选供体投影图}
\end{figure}
<% if fmt_plot_interpretation %>
\begin{infobox}{图形解读}
<{ fmt_plot_interpretation }>
\end{infobox}
<% endif %>
<% endif %>
<% elif fmt_degraded %>
\begin{warnbox}{FMT 配型状态}
<{ fmt_message }>
\end{warnbox}
<% elif fmt_unavailable %>
\begin{warnbox}{FMT 配型状态}
<{ fmt_message }>
\end{warnbox}
<% endif %>
\begin{infobox}{FMT 配型说明}
<{ fmt_explanation }>
\end{infobox}
<% if fmt_patient_interpretation %>
\begin{infobox}{患者解读}
<{ fmt_patient_interpretation }>
\end{infobox}
<% endif %>
<% endif %>

<% if llm_advice %>
\begin{recommendbox}{个体化建议}
<{ llm_advice }>
\end{recommendbox}
<% endif %>

\section{参考资料}
<% if references %>
\begin{enumerate}
<% for ref in references %>
  \item <{ ref }>
<% endfor %>
\end{enumerate}
<% else %>
\begin{infobox}{参考资料状态}
<{ references_empty_message }>
\end{infobox}
<% endif %>

<% endblock %>
